\section{Einleitung}
Die nachfolgende Arbeit beschäftigt sich mit der neuronalen Klassifikation dank gelernter Konzepte, die sich in Netzwerken wiederfinden.
Hierzu wurde im  ``Independent Coursework I'' zuvor ein ``Pruning eines neuronalen Netzwerks anhand von KI-Erklärbarkeit'' ausgearbeitet. 
Anhand der über die Layerwise Relevance Propagation (LRP) \cite{DBLP:journals/corr/abs-1912-08881} ermittelten Relevanzen, die die Erklärbarkeit von Netzwerkentscheidungen transparenter machen, konnte das Netzwerk um Einiges verkleinert werden. 
Dagegen die durch das ProtoPNet \cite{Mohammadjafari2021Using} erreichte Erklärbarkeit sorgte für das Offenlegen von Prototypen, die aussagekräftig zur Klassifizierung beitragen. Diesbezüglich konnte das Pruning dieser für eine genauere Distinktion zwischen den Klassen hinreichend sein. 

Auf diesem Hintergrund und in Anlehnung an das ``Forschungsprojekt B'' werden in dieser Arbeit dem Forschungsprojekt ähnliche Rahmenbedingungen geschaffen: Anhand des am CBMI \cite{cbmi} enstandenen Mitosen-Klassifikators\footnote[1]{Hierzu werden ebenso das zugehörige \pythoninline{state_dict_path}, der Datensatz \pythoninline{examples_ds_from_MP.train.HTW.train} und der zur Verfügung gestellte Modell-Quellcode genutzt} besteht der Fokus auf der Untersuchung der Klassifikations-Konzepten, die näher inspiziert werden können.


In Anbetracht dessen wird auf dem in dem ``Forschungsprojekt B'' vorgestellten Ansatz zum Auffinden genutzt: Die bereits vorgestellte Concept Relevance Propagation (CRP) \cite{achtibat2022towards} eignet sich, um relevante und irrelevante Konzepte hinsichtlich einer Klassifikation zu identifizieren. Weiter kann mittels dem Python-Package Zennit-CRP \cite{achtibat2022crpzennit} dieses Vorhaben praktisch umgesetzt werden.

Im Folgenden wird daher untersucht, wie diese erlernten Konzepte selektiv zu einer Mitosen-Klassifikation beitragen - und wie sich diese nach vorherigem, dem Forschungsprojekt entstandenem Pruning verändert.
